\cotefile{34}
\scanpage{67}
\mspage{II 15}
% archive sheet number: 66
% lines 1–4
\deleted{\unsure{choisissons} $K'$ assez grand, et utilisons le cor. 2 précédent, \unsure{pour} \deleted{\unsure{donc}} \unsure{supposer} que l'action de $\operatorname{Gal}(\overline{K'}/K')$ est \unsure{triviale}. Changeant de notations, on peut supposer alors $K'=K$. Or l'action \unsure{en est triviale}.}

% lines 5–6
Il suffit d'établir le résultat suivant (qui a été signalé par Serre) :

% lines 7–13
\uline{Lemme} Soit $V$ anneau de val. discrète, corps fractions $K$, $A$ schéma abélien sur $K$, \inserted{$\ell$ un nb premier distinct de la car. de \uline{$K$}, (mais pas néc. de la car. du corps résiduel)} supposons que l'action du groupe d'inertie \deleted{\unsure{sur}} $I\subset\operatorname{Gal}(\bar K/K)$ sur $E=T_\ell(A(\bar K))\otimes_{\unsure{\mathbb Z_\ell|Z_\ell}}\unsure{\mathbb Q_\ell|Q_\ell}$ soit \deleted{\unsure{triviale}} \deleted{\unsure{abélienne}} demi-simple, \deleted{\unsure{(dans ce cas abélienne, alors}} alors \deleted{$A$ a bonne} (et alors seulement) $A$ a bonne réduction potentielle''.

% line 13, left margin
% slanted
\margin{\uline{Pour la démonstration} \uline{\unsure{renvoyons|renvoyer} à exposé} \uline{ultérieur}}

% lines 14–17
Le cas de la car. résiduelle $\ne\ell$ est évident, car l'action de $I$ étant un. unipotente (\unsure{J|I}), \inserted{est semi-simple, ssi elle} doit être un. triviale. \deleted{\unsure{et}} Le cas de la car résiduelle $\ell$ (donc $K$ de car zéro) demande un traitement spécial.

% lines 18–24
On peut supposer $V$ \deleted{hensélien, et \uline{complet}, et, après \unsure{ext}} \inserted{à \uline{corps résiduel sép. clos}.} Considérons le modèle de Néron \deleted{\unsure{de}} $G$ de $A$ sur $V$ [\quad]. \deleted{Alors} Le th. de réduction \unsure{$\frac12$ st. R.}, qui sera prouvé ultérieurement, indique que (quitte à faire une extension finie sur $K$) on peut supposer que $G_{\unsure{R|\kappa|k}}^0$ est ext. d'une V.A. par un tore. \inserted{$T$ de dimension $r$} On sait donc que $T_\ell(A(\bar K))$ est extension d'un g.
\scanpage{68}
\mspage{II.16}
% archive sheet number: 67
% lines 1–9
\deleted{groupe \unsure{à}} \inserted{$\unsure{\mathbb Z_\ell|Z_\ell}$ module $M\simeq$} $(\deleted{\unsure{\mathbb Q_\ell}}\mathbb Z_\ell)^r$ \uline{sur lequel} $I$ \uline{opère trivialement}, par un \deleted{groupe} $N$ abélien de rang $2g-r$, \inserted{can. isom. à $T_\ell(G(\unsure{S|s}))$} Comme l'action de $I$ est semi-simple, il s'ensuit que l'on a un ss-module $M\simeq\mathbb Z_\ell^r\subset T_\ell(A(\bar K))$ \uline{sur lequel} $I$ \uline{opère trivialement}, \deleted{\unsure{donc}} \inserted{avec $M\cap N=0$. On en déduit} un sous-groupe de la forme $(\unsure{\mathbb Q_\ell|Q_\ell}/\mathbb Z_\ell)^r$ du grp $A(K)$ : Par définition du modèle de Néron, il sera contenu dans $G(S)$, i.e. $M\subset N$, donc $M=0$, donc $r=0$, ce qui signifie que la réduction est \unsure{st.R|b.R}, cqfd.

\par\noindent\rule{.45\linewidth}{.4pt}\par\noindent\rule{.42\linewidth}{.4pt}\par

% lines 10–14
\uline{Remarque de Serre} Si on savait montrer que (dans le cas de la car. nulle) action résoluble implique action ess. abélienne, \inserted{(VA sur corps de type fini)} on pourrait en conclure qu'une \inserted{telle} VA est ``de type CM'' après extension finie du corps de base.

% lines 15–21
On est ramené, par ce qui précède, au cas d'un \deleted{VA} \inserted{schéma abélien} sur l'anneau des entiers \inserted{$R$} d'un corps de nbs $K$, avec action abélienne de $\operatorname{Gal}(\bar K/K)$ donc de $\pi_1(S,\unsure{\xi|\eta})$ ($S=\operatorname{Spec}R$). L'action de $\pi_1(S,\xi)$ \inserted{sur $E_\ell=T_\ell(A(\bar K))_{\mathbb Q_\ell}$} est loc algébrique \inserted{en tout $m\in\operatorname{Max}(R)$} (conditions locales sur l'action des $\pi_1$ des corps locaux, i.e. des groupes de décomposition) et est ``\inserted{loc} rationnelle'', \inserted{i.e. les pol. car. des \unsure{(v.b. de Frob)|subst. de Frob} à coeff $\in\mathbb Q$.} \deleted{\unsure{car}} Or on prouve que si on a deux telles représentations, \deleted{\unsure{alors}} un $E_\ell$ \uline{et}
\scanpage{69}
\mspage{II.17}
% archive sheet number: 68
% lines 1–4
dans un $E_{\ell'}$, l'une étant semi-simple abélienne, \deleted{\unsure{donc}} \deleted{l'autre est} avec les $\unsure{\rightrightarrows|=}$ pol. car. de Frob, pour presque tous $m\in\operatorname{Spec}R$, alors l'autre représentation est \uline{résoluble},
% lines 4–7, cancelled block including interlinear revisions
\deleted{\unsure{qui} \unsure{a} \unsure{pour} \unsure{obj.} $E_\ell,E_{\ell'}$ \unsure{provenant} $A$ \unsure{par} $T_\ell,T_{\ell'}$, Supposons que \deleted{\unsure{elle}} \unsure{donc} qu'elle soit \inserted{ess. résoluble} \deleted{\unsure{ess.}} aussi abélienne. Alors \unsure{on trouve} que \inserted{et la représentation $\frac12$ simple associée} \unsure{et elle donne} la représentation \uline{sur $\unsure{\mathbb Q|Q}$}}
% lines 8–10
et pour tt $m\in\operatorname{Spec}R$, la repr. \uuline{$\frac12$ simple} \uuline{associée} à $E_{\ell'}$ \unsure{ou|on} \deleted{\unsure{elle|est}} \deleted{\unsure{alors|elle}} est \unsure{le|la} $\unsure{=|\text{même}}$ que celle qui est définie par $E_\ell$. Choisissons alors $\ell$ tel que
\scanpage{70}
\mspage{II.18}
% archive sheet number: 69
% lines 1–3
\deleted{pourrait généraliser substantiellement \deleted{\unsure{le}} 1.3, 1.5, \deleted{\illegible{1}}, \deleted{\illegible{1}} 1.9 \unsure{$\#$} (ss-réserve de $\frac12$ simplicité), 1.11, en y remplaçant partout ``ess. abélienne'' par ``ess.$^{\mathrm t}$ résoluble''.}

% line 4
2. \uline{Démonstration} \inserted{de 1.1} : \uline{Cas où $X$ est de dim} 1.

% lines 5–9
\deleted{Alors $X$ est un ouvert}

2.1. Notons d'abord (sans hyp. sur la dim de $X$) qu'on peut supposer $X$ \uline{affine}. \deleted{\unsure{du schéma géom. irréduct.}} En effet, si $U$ est un ouvert affine \deleted{\unsure{de}} de $X$ contenant l'image de $\unsure{\xi|\eta}$, on a \deleted{On sait que $\pi_1($} \inserted{un \unsure{hom.} de suites exactes}
% lines 10–12
\diagram{\begin{tikzcd}[ampersand replacement=\&]
1 \arrow[r,dash] \& \pi_1(\bar U,\xi) \arrow[r,dash] \arrow[d,"\alpha"] \& \pi_1(\deleted{\unsure{X}}U,\xi) \arrow[r,dash] \arrow[d,"\beta"] \& \pi_1(S',\xi) \arrow[r,dash] \arrow[d,Rightarrow] \& 1 \\
1 \arrow[r,dash] \& \pi_1(\bar X,\xi) \arrow[r,dash] \& \pi_1(X,\xi) \arrow[r,dash] \& \pi_1(S,\xi) \arrow[r,dash] \& 1 .
\end{tikzcd}}

% lines 13–17
Comme \deleted{$X$ est} $\bar X$ \deleted{\unsure{est}} \deleted{géométriquement \unsure{unibranche}} \inserted{normal}, \deleted{(\unsure{ainsi} $\bar X$ cf. EGA IV 6.15.6)}, il s'ensuit que $\beta$ est surjectif (SGA1 V \unsure{\ldots}), donc d'après \deleted{$\pi_1(X$ $\alpha$ et $\beta$ sont surjectifs.} On en conclut que si la \deleted{\unsure{même}} conclusion de 1.1. est vraie pour $(U,\xi)$, elle est vraie pour $(X,\xi)$.

% lines left margin beside 17–24
% margin note written sideways
% slanted
\margin{\deleted{2.2. On voit de même que la conclusion \unsure{ne change pas} si on \unsure{remplace} $K$ par une \unsure{ext.} finie $K'$, et $X$ par $X\otimes_K K'$, et $\xi$ par $\xi'$ \unsure{au-dessus de} $\xi$ : cela revient à \unsure{remplacer} $\pi_1(X,\xi)$ par un ss-groupe ouvert, \unsure{l'image réciproque d'un} ss-groupe ouvert de $\pi_1(S,\xi)$.}}

% lines 18–24
\deleted{2.3. \deleted{\unsure{Soit donc}} Supposons donc \deleted{\unsure{X}} $X$ affine, donc quasi-projectif, et supposons de plus $\dim X=1$. Alors $X$ est un \deleted{\unsure{normal|sous-sch}} ouvert d'une courbe normale projective $\hat X$ (EGA \unsure{II|III} \unsure{7.1|7.}). Quitte à passer à une extension \deleted{radicielle} finie de $K$, on est ramené au cas où $\hat X'$ est \uline{lisse} sur $K$, et où les points de \deleted{\unsure{X}} $\hat X-X$ sont tous rationnels sur $K$. Soit $x_i$ ($1\leq i\leq n$) ces pts.}
\scanpage{71}
\mspage{II.19}
% archive sheet number: 70
% line 1
\deleted{On a alors la suite exacte bien connue :}

% lines 2–7
2.2. On voit de $\unsure{=|\text{même}}$ que la conclusion ne change pas \deleted{$\unsure{\pi_1(X}$} si on remplace $\unsure{k|h}$ par une extension finie $k'$, $X$ par $X'=X\otimes_k k'$, et $\unsure{\xi|\eta}$ par un pt géom $\xi'$ au-dessus de $\xi$. En effet, cela revient à \deleted{\unsure{en}} remplacer $\pi_1(X,\xi)$ par un ss-groupe ouvert, image inverse du ss-groupe ouvert de $\pi_1(S,\xi)$ défini par l'extension $k'/k$ de $k$.

% lines 8–12
2.3. Supposons maintenant $X$ quasi-projectif, comme il est loisible par 2.1. Il existe alors un schéma \deleted{\unsure{quasi}} projectif $X'$ sur $k$, \deleted{\unsure{un}} tel que $X$ soit isom à un ouvert \deleted{\unsure{de}} \inserted{dense} de $X'$. On peut de plus supposer $X'$ normal (quitte à le remplacer par \inserted{le normalisé de} $X'_{\mathrm{red}}$).
% lines 13–18
Utilisant \deleted{\unsure{le}} EGA IV \unsure{17.15.14.2}, on voit qu'il \deleted{\unsure{existe}} existe une extension finie radicielle $k_1$ de $k$, telle que le normalisé \inserted{$\unsure{\tilde X'_1|\hat X'_1}$} de $X'_1\deleted{\unsure{=} }=X'\otimes_k k_1$ soit géom\textsuperscript{t} normal. \deleted{\unsure{Il en est donc}} Il en est de $=$ du normalisé \inserted{$\hat X_1$} de $X_1=X\otimes_k k_1$, qui s'identifie à \deleted{\unsure{de}} un ouvert de $\tilde X'_1$. \deleted{\unsure{il est}}
% lines 19–24
Comme $X_1$ est géom\textsuperscript{t} unibranche (EGA IV 6.15.6), $\unsure{\tilde X_1|\tilde X'_1}\unsure{-|\longrightarrow}X_1$ est radiciel surjectif, donc induit un \uline{isom} sur les $\pi_1$ (SGA 1 \uuline{VIII} \unsure{\ldots}), \inserted{et de $=$ pour $\widetilde{\bar X}_1-\bar X_1$.} Donc, conjuguant cette remarque avec 2.2, on voit qu'on est ramené au cas où $X'$ est géom\textsuperscript{t} normal.
\scanpage{72}
\mspage{II.20}
% archive sheet number: 71
% line 1; insertion at loop, written in box at foot
2.4. \inserted{faire d'abord une réduction au cas où $\pi=\pi_1(X,\unsure{\xi|\eta})$, en introduisant \deleted{\unsure{sur}} le rev. étale fini $X'$ de $X$ qui correspond à $\pi\subset\pi_1(X,\xi)$.}
% lines 1–6
On sait déjà que $V=\pi_1(\bar X)^{\mathrm{ab}}(\ell)$ est un module de t.f. sur $\unsure{\mathbb Z_\ell|Z_\ell}$, \deleted{ce qui} ce qui revient au $\unsure{=|\text{même}}$, que $H^1(\bar X,\mathbb Z/\ell\mathbb Z)$ est un groupe fini. \deleted{\unsure{(}} Ceci est vrai pour tt schéma de t.f. sur un corps \deleted{alg.} \inserted{sép.} clos, et se \inserted{peut} démontrer par réduction au cas des courbes, par le procédé du n° 3.
% lines 7–9
Lorsqu'on \deleted{\unsure{admet}} \inserted{sait} que $\pi_1(\bar X)^{\mathrm{ab}}(\ell)$ est de t.f. sur $\mathbb Z_\ell$, alors la relation à prouver, \deleted{\unsure{qui s'écrit}} s'écrit aussi \deleted{$\pi_1(\bar X)(\ell)_{\pi_0}$ fini i.e.}
% line 10
\[
\deleted{\unsure{(2.4)}}\qquad V_{\pi_0}\ \text{fini}\quad\text{i.e.}\quad (V\otimes_{\mathbb Z_\ell}\unsure{\mathbb Q_\ell|Q_\ell})_{\pi_0}=0,
\]
% line 11
et par dualité elle équivaut à la relation
% line 12
\[
(\check V)^{\pi_0}=0
\]
% lines 13–15
où $\check V=\operatorname{Hom}_{\mathbb Z_\ell}(\unsure{W|V|N},\mathbb Z_\ell)$. Or on a un isom canonique
\[
\check V\simeq\operatorname{Hom}(\pi_1(\bar X,\xi),\mathbb Z_\ell)\simeq H^1(\bar X,\mathbb Z_\ell),
\]
% lines 16–18
évidemment compatible avec l'action de $\pi_0$. Donc la conclusion de 1.1 s'écrit de façon équivalente, sous forme cohomologique
% line 19
\[
H_{\deleted{\unsure{1}}}^1(\bar X,\mathbb Z_\ell)^{\unsure{\bar\pi_0|\pi_0}}=0.\tag*{(2.4.1)}
\]
% lines lower left margin
\margin{faire une section 2.4\unsure{\ldots} spéciale pour cette \unsure{réduction}, \deleted{\unsure{pourrait déjà}} \uline{être \unsure{venu|revu}\unsure{e}}.}
\scanpage{73}
\mspage{II.21}
% archive sheet number: 72
% lines 1–5
2.5. Dans ce qui suit, nous supposerons $\dim X=1$, donc $\dim X'=1$. Quitte à faire encore une extension finie sur $\unsure{k|h}$, on peut supposer les points $x_i$ ($1\leq i\leq n$) de $X'-X$ rationnels sur $k$. On a alors une suite exacte bien connue
% lines 6–7
\noindent (2.5.1)
\diagram{\begin{tikzcd}[ampersand replacement=\&,column sep=small]
\unsure{0|1} \arrow[r,dash] \& H^1(\bar X',\unsure{\mathbb Z_\ell|Z_\ell}) \arrow[r] \& H^1(\bar X,\mathbb Z_\ell) \arrow[r] \& \displaystyle\prod_{\deleted{\unsure{0}}1\leq i\leq n}H_x^2(\bar X,\mathbb Z_\ell) \arrow[r] \& H^2(\bar X',\mathbb Z_\ell) \arrow[d] \\
\&\&\&\&0
\end{tikzcd}}
% lines 8–10
\deleted{qui est} \inserted{dont les flèches sont} évidemment compatibles avec l'action de \deleted{$\unsure{\pi_0}$} \deleted{$\unsure{k}$} $\pi_0=\operatorname{Gal}(\bar k/k)$. De plus, on a un isomorphisme canonique
% line 11
\[
H_x^2(\bar X,\mathbb Z_\ell)\simeq\mathbb Z_\ell(-1)\qquad \deleted{\unsure{(\varprojlim_\nu{}_{\ell^\nu}\bar k^*)}}\tag*{(2.5.2)}
\]
% line 12
où on a
% line 13
\[
\deleted{\unsure{H^1}}\mathbb Z_\ell(-1)=\mathbb Z_\ell(1)^{\otimes(-1)},\qquad \mathbb Z_\ell(1)=\varprojlim_{\unsure{\nu|v}}({}_{\ell^\nu}\bar k^{\unsure{*|\times}}).\tag*{(2.5.3)}
\]
% lines 14–19
(De plus, on a un isom. canonique $H^2(\bar X',\mathbb Z_\ell)\simeq\mathbb Z_\ell(-1)$, et l'hom $\prod\longrightarrow H^{\unsure{1|2}}(\bar X',\mathbb Z_\ell)$ s'identifie à l'hom somme). L'isomorphisme (2.5.2) est encore compatible avec l'action de $\pi_0$, opérant sur le deuxième m\textsuperscript{bre} de la façon explicitée par la formule (2.5.3). \deleted{\unsure{Il s'}}

% lines 20–23
Il s'ensuit qu'en tant que $\overset{\deleted{\unsure{-}}}{\pi}_0$-module $\ell$-adique, $H^1(\bar X,\mathbb Z_\ell)$ est une extension de $\mathbb Z_\ell(1)^{n-1}$ par $H^1(\bar X',\mathbb Z_\ell)$. Comme l'opération de $\pi_0$ sur $\mathbb Z_\ell(1)$ est non triviale, il s'ensuit \inserted{déjà} qu'on a
\scanpage{74}
\mspage{II.22}
% archive sheet number: 73
% line 1
\[
H^1(\bar X,\unsure{\mathbb Z_\ell|Z_\ell})^{\pi_0}\subset H^1(\bar X',\mathbb Z_\ell)^{\pi_0}.\tag*{(2.5.4).}
\]
% lines 2–3
\deleted{Il reste à \deleted{\unsure{Démontrer}} \unsure{voir}} que le dernier membre est nul :
\[
H^1(\bar X',\mathbb Z_\ell)^{\pi_0}=1.\tag*{(2.5.5)}
\]
% lines 4–9
C'est une conséquence bien connue de l'hyp. de Riemann-Weil, prouvée en dim. 1. Explicitons l'argument. \deleted{Soit $A$ \deleted{\unsure{la}} \unsure{la jacobienne},} \inserted{Posons $X'=C_k$.} On sait que $\unsure{k|h}$ est le \deleted{\unsure{anneau}} \inserted{corps} des fractions d'une \deleted{\unsure{algèbre}} Algèbre \inserted{$R$} de t.f. sur $\mathbb Z$, \deleted{\unsure{donc}} telle que \deleted{\unsure{X}} $C_k$ provienne d'un schéma $C_S$ sur $S=\operatorname{Spec}R$, propre et lisse sur $S$. \inserted{On peut supposer $S$ \uline{normal}.} La \deleted{homomorphisme} \inserted{relation} (2.5.5) s'écrit alors
% line 10
\[
H^0(S,R^1f_*(\mathbb Z_\ell))=0
\]
% lines 11–16
$f:C-S$. Sous cette forme, on voit que cette relation \deleted{\unsure{entraîne}} ne change pas si on \deleted{change} \inserted{remplace} le pt géom $\unsure{\xi|\eta}$ de $S$ par \deleted{\unsure{parallèle}} \inserted{un} autre $\xi'$, \inserted{en faisant opérer $\pi_1(S,\xi')$ sur $H^1(C_{\xi'},\mathbb Z_\ell)$,} \deleted{\unsure{qu'on}} peut supposer \inserted{\unsure{le choisir en un|le choix en un}} \deleted{fini} \uline{fermé} \inserted{$s$} de $S$, donc à corps résiduel fini. \deleted{Alors} À fortiori, il suffit de prouver la relation
% line 17
\[
H^1(C_{\bar s},\mathbb Z_\ell)^{\pi_1(s,\bar s)}=0\tag*{(2.5.6)}
\]
% lines 18–22
Cela nous ramène au cas d'un corps de base \uline{fini}. Mais dans ce cas, on sait \inserted{$\pi_0\simeq\hat{\mathbb Z}$,} \deleted{\unsure{d'après les}} \inserted{et à l'effet \unsure{le|de} prouver, il suffit de voir que \unsure{Ker} $(1-f)$ dans $H^1(C_{\bar s},\mathbb Z_\ell)\simeq T_\ell(A)^\vee$ est nul, i.e Coker $(1-f)$ dans $T_\ell(A)$ fini, i.e. Ker $(1-f)$ dans $T_\ell(A)$ fini, i.e. car $A(k)$ fini, ce qui est trivial.} \deleted{les valeurs propres du générateur de Frobenius, frob de $\hat{\mathbb Z}$ opérant sur $H^1(C_{\bar s},\mathbb Z_\ell)$ sont de valeur absolue $q^{1/2}$ ($q=\operatorname{Card}k$). \inserted{[2,\unsure{\ldots}]} Donc elles sont $\ne1$, \unsure{d'où} (2.5.6).}

% lines left margin beside 8–9, circled and crossed
\margin{\deleted{\unsure{f fait un} \illegible{2} \unsure{pt.} EGA IV \unsure{\ldots}}}
% lines left margin beside 11–19
% margin note written sideways
\margin{\uline{Remarque} 2.6. Nous avons utilisé exactement 2 propriétés de notre corps $K$ : a) \unsure{l'extension} $\subset\bar K$ obtenue en adjoignant les racines $\ell^{\unsure{\nu|v}}$-ièmes de l'unité, pour $\forall\nu$, est infinie b) Pour tout schéma abélien $A$ sur une extension finie $K'$ de $K$, $T_\ell(A)_\pi$ est fini, où $\pi=\operatorname{Gal}(\overline{K'}/K')$. Pour b) il \uline{ne suffit pas} de savoir que $T_\ell(A)^\pi$ est tjrs nul (exemple : $K=\unsure{\mathbb Q_\ell|Q_\ell}$)!}
\scanpage{75}
\mspage{II.23}
% archive sheet number: 74
% line 1
3. \uline{Réduction au cas de la dimension} 1.

% lines 2–6
3.1. \deleted{\unsure{Supp}} Soit d'abord $\unsure{k|h}$ un \deleted{\unsure{groupe}} corps quelconque, $X'$ un schéma \deleted{géométriquement irréductible} \inserted{\deleted{\unsure{de}} projectif \inserted{géom} normal, géom connexe,} sur $k$, $X$ un ouvert dense de $X$, \deleted{On se propose d'} \inserted{$\unsure{\xi|\eta}$ un pt géom. de $X'$, $\bar k$ la clôture séparable de $k$ dans $\kappa(\xi)$, $\bar X',\bar X$ déduits de $X',X$ par $k\unsure{-|\longrightarrow}\bar k$.} On veut étudier $\pi_1(\bar X^{\unsure{\sharp|*|\prime}},\xi)$, et \deleted{\unsure{de}} le majorer en fonction du groupe fondamental d'une \uline{courbe} lisse sur un corps convenable.
% lines 7–11
Comme on veut tenir compte des opérations, possibles de $\pi_1(S,\xi)$ ($S=\operatorname{Spec}k$) sur $\pi_1(\bar X^\sharp,\xi)$ resp sur sa partie abélienne, on va considérer aussi la suite exacte
% line 12
\noindent (3.1.1)
\diagram{\begin{tikzcd}[ampersand replacement=\&]
1 \arrow[r,dash] \& \pi_1(\bar X,\xi) \arrow[r,dash] \& \pi_1(X,\xi) \arrow[r,dash] \& \pi_1(S,\xi) \arrow[r,dash] \& 1 .
\end{tikzcd}}

% lines 13–19
\deleted{Soit $T$ le schéma qui} Nous supposons \deleted{que} $X'$ immergé dans un $\unsure{\mathbb P^r_k|P^r_k}$. \deleted{$\unsure{P^r_k}$ \unsure{Soit}} Soit $n=\dim X'$, et soit $T$ le grassmannien des sous-espaces projectifs de codimension $n-1$ \deleted{\unsure{de}} dans $P$, \inserted{(EGA I, $2^{\text{ième}}$ édition\unsure{\ldots})} $t$ le point générique de $T$, $K=k(t)$, qui est une extension \deleted{pure} \inserted{de t.f.} de $k$. \deleted{Alors} Soit $L_t$ le sous-espace \inserted{projectif} de codim $n-1$ \inserted{``générique''} de $P_K$ qui correspond à $t\in T_K$, et considérons $C'=L\cap X'_K$, \deleted{qui est} Il est connu que $\dim C'=1$, que $C'$ est géom. normal et géom connexe. Soit $C$ l'image inverse de $X_K$ dans $C'$.
% lines left of 18–21
\diagram{\begin{tikzcd}[ampersand replacement=\&]
C \arrow[d] \& C'=L_t\cap X'_K \arrow[l] \arrow[d,dash] \\
X_K \arrow[r,phantom,"\subset"] \& X'_K
\end{tikzcd}}
\scanpage{76}
\mspage{II. 24}
% archive sheet number: 75
% lines 1–3
Considérons \deleted{\unsure{l'image de $C$}} le morphisme canonique
\[C\longrightarrow X\]
composé de l'inclusion $C\hookrightarrow X_K$ et $X_K-X$.
% lines 4–9
Nous supposons que \deleted{\unsure{le}} l'image\inserted{$x$} de $\xi$ dans $X$ est contenue dans celle de $C$. Ce qui signifie $\dim\overline{\{x\}}\geq1$); \deleted{\unsure{il en résulte}} comme l'extension 3.1.1. ne dépend pas, à isom. près, du choix de $\xi$, cela est une hypothèse \unsure{anodine|admissible}. Soit $\xi_1$ un point géométrique de $C$ au dessus de $\xi$
% lines 10–13
\[(3.1.2)\]
\diagram{\begin{tikzcd}[ampersand replacement=\&]
\xi_1 \arrow[r,dash] \arrow[d,dash] \& \xi \arrow[d]\\
C \arrow[r,dash] \arrow[d] \& X \arrow[d,dash]\\
\operatorname{Spec}K \arrow[r,dash] \& S
\end{tikzcd}} $*$
et soit $\bar K$ la clôture sép. de $K=k(t)$ dans $\kappa(\xi_1)$.
% lines 14–18
\deleted{On en conclut un diagramme \inserted{commutatif} d'inclusions de corps}
\diagram{\begin{tikzcd}[ampersand replacement=\&]
\deleted{k} \arrow[r,dash] \arrow[d,dash] \& \deleted{K=k(t)} \arrow[dd,dash]\\
\deleted{\bar k} \arrow[dr] \& \\
\& \deleted{\bar K}
\end{tikzcd}}
% lines 19–22
\deleted{Considérons aussi le composé $K\bar k=\unsure{L|l}$ dans $\bar K$, qui n'est autre que le corps $\kappa(\bar t)$, où $\bar t$ est le point générique de $T_{\bar k}$.}
% lines 23–25
Le diagramme (3.1.2) donne un diagramme commutatif, à lignes exactes :
\scanpage{77}
\mspage{II. 25}
% archive sheet number: 76
% lines 1–4
\[(3.1.3)\]
\diagram{\begin{tikzcd}[ampersand replacement=\&,column sep=small]
1 \arrow[r,dash] \& \pi_1(\bar C,\xi_1) \arrow[r] \arrow[d,"\alpha"'] \& \pi_1(C,\xi_1) \arrow[r,dash] \arrow[d,"\beta"] \& \begin{gathered}\inserted{\operatorname{Gal}(\bar K/K)}\\\deleted{\unsure{\operatorname{Gal}(\bar K/K)}}\\\deleted{\unsure{\pi_1(\operatorname{Spec}K,\xi_1)}}\\\pi_1(\operatorname{Spec}K,\xi_1)\unsure{=}\end{gathered} \arrow[r,dash] \arrow[d,"\gamma"] \& 1\\
1 \arrow[r,dash] \& \pi_1(\bar X,\xi) \arrow[r,dash] \& \pi_1(X,\xi) \arrow[r,dash] \& \pi_1(\operatorname{Spec}k,\xi) \arrow[r,dash] \& 1
\end{tikzcd}}
,
% lines 5–6
où on a posé $\bar C=C\otimes_K\bar K$. \deleted{\unsure{Ceci prouve que}} Ceci précise le
% lines 7–9
\uline{Lemme} 3.2. Avec les \inserted{hyp. et} notations précédentes, dans le diagramme (3.1.3), les \inserted{trois} hom \deleted{$\gamma$ et} \inserted{verticaux} \deleted{\unsure{$\alpha$ sont des}} sont \uline{surjectifs}.
% lines 10–15
\deleted{\unsure{Dém.} Pour prouver le lemme}, Il suffit de prouver la surjectivité de $\gamma$ et $\alpha$. Celle de $\gamma$ provient du fait que $K=\unsure{k(t)}$ est une extension \deleted{transcendante} pure de $k$, donc \deleted{\unsure{le morphisme}} \deleted{\unsure{de}} $k$ est sép. fermé dans $K$.
% lines 16–23
Notons maintenant que la flèche $\bar C\longrightarrow\bar X$ donnant naissance à $\alpha$ est la même que celle qui \unsure{s'introduit|s'obtient} en \deleted{\unsure{remplaçant le cas}} faisant le changement de corps de base $k-\bar k$. Donc on peut supposer $k$ séparablement clos, et \deleted{\unsure{il est}} on \deleted{est ramené à \unsure{pr}} \inserted{doit examiner le morphisme sur les $\pi_1$ \uline{induit} par}
\[\unsure{\bar C|\underline C}-X.\]
% lines 24–26
Dire que le morphisme sur les $\pi_1$ est surjectif signifie que pour tt revêtement étale connexe $\bar X_1$ de $X$, son image inverse $\bar C_1$ sur $\bar C$ est
\scanpage{78}
\mspage{II. 26}
% archive sheet number: 77
% lines 1–7
connexe, \inserted{\unsure{i.e.} son image inverse $\bar C_1$ sur $C$ est géom. connexe, ou même géom. irréduct.} Or $X_1$ provient d'un revêtement \deleted{étale} \inserted{normal} intègre \deleted{\unsure{ramifié}} \inserted{[EGA \unsure{V|IV}]} $X'_1-X'$. En vertu du th. de Bertini, l'image inverse \deleted{\unsure{de}} dans $X'_1\longrightarrow\mathbb P^r$ \deleted{\unsure{est géom}} de $L_t$ est \deleted{géom.} géom. irréductible, ce qui implique que l'ouvert $C_1^{\unsure{'|-}}$ tracé sur $C$ l'est également, i.e. $\bar C_1$ est géom connexe. Cela prouve le lemme.
% lines 8–14
3.3. De 3.2 résulte aussitôt que si \deleted{l'extension} $C$ sur $K$ satisfait aux conclusions de 1.1, il en est de même de $X$ sur $k$. Compte tenu de la réduction déjà faite dans 2.3, cela nous ramène, pour prouver 1.1 \inserted{sur $k$}, au cas d'une courbe algébrique \inserted{sur $K$}, déjà traité au n° 2.

(NB. Si $k$ est de t.f., il en est de même de $K$ !).
% lines 15–23
\uline{Remarque} 3.4. Le lemme 3.2. permet aussi de montrer que si $X$ est un schéma de t.f. sur un corps séparablement clos \inserted{de car. $p>0$}, alors le \deleted{\unsure{petit}} plus grand \inserted{groupe} quotient profini premier à $p$ de $\pi_1(X,\xi)$, \inserted{soit $\pi_1(X,\xi)'$} est topologique à engendrement fini. En effet, la théorie de la descente \inserted{(pour les groupes fond$^{\unsure{aux}}$ (SGA 1 \unsure{IX|VIII}))} nous permet de nous ramener au cas où $X$ est géom. unibranche irréductible, \deleted{\unsure{et}} puis au cas \deleted{\unsure{d'un}} $X$ affine, puis 3.2. nous ramène au cas où $X$ \deleted{\unsure{d'une}} \inserted{est un ouvert} d'une courbe projective lisse $X'$. Or pour un tel $X$, cela résulte de la structure connue de $\pi_1(X,\xi)'$ (SGA 1\ldots).
\scanpage{79}
\mspage{II. 27}
% archive sheet number: 78
% line 1
4. \uline{Un complément au th. principal}.
% lines 2–6
\uline{Proposition} 4.1 Sous les conditions de 1.1, pour tout \deleted{\unsure{$\ell$}} nb premier $\ell$ sauf un nb fini, \deleted{le groupe} \deleted{\unsure{$\pi$}} on a $\overline{[\pi',I(\ell)]}=I(\ell)$, i.e. le plus grand quotient \unsure{profini|propre} de $I(\ell)$ sur lequel $\pi_0$ opère trivialement est \uline{nul}.
% lines 7–10
On voit par dualité, en procédant comme dans 2.4, que cet énoncé équivaut au suivant :

\deleted{\uline{Proposition}} \inserted{Corollaire} 4.2. Pour presque tout $\ell$, on a
\[H^1(\bar X,\mathbb Z/\ell\mathbb Z)^{\pi_0}=0.\]
% lines 11–17
Pour prouver ceci, on est encore ramené, \deleted{\unsure{comme}} \deleted{\unsure{au cas où}} \deleted{\unsure{qu'on a}} par la réduction de 2.3, au cas où \deleted{$X\subset X'$} $X$ est un ouvert d'un $X'$ projectif et géom normal sur $k$. Dans ce cas, \deleted{le lemme} 3.2. nous ramène au cas où de plus $X'$ est de dim 1, donc est une courbe lisse et complète.
% lines 18–21
Quitte à \deleted{\unsure{agrandir $k$}} passer à une extension finie de $k$, on peut encore supposer les pts \inserted{$x_i$} de $X'-X$ rationnels sur $k$. On utilise alors la suite exacte analogue à (2.5.1)
% lines 22–24
\[(4.2.1)\]
\diagram{\begin{tikzcd}[ampersand replacement=\&,column sep=small]
{} \arrow[r,dash] \& H^1(\bar X',\mathbb Z/\ell\mathbb Z) \arrow[r] \& H^1(\bar X,\mathbb Z/\ell\mathbb Z) \arrow[r] \& \deleted{\unsure{H}}\prod_i H^2_x(\bar X,\mathbb Z/\ell\mathbb Z) \arrow[r,dash] \& H^2(\bar X',\unsure{\mathbb Z(\ell)|\mathbb Z/\ell\mathbb Z}) \arrow[d,dash]\\
\&\&\&\&{}
\end{tikzcd}}
% lines 25–28
qui nous montre que $H^1(\bar X,\mathbb Z/\ell\mathbb Z)$ est une extension de $\bigl(\mu_\ell(\bar k^*)^{\otimes-1}\bigr)^{n-1}$ par $H^1(\bar X',\mathbb Z/\ell\mathbb Z)$.

\deleted{\unsure{Comme le gr}} Soit
\[N=\text{card du grp des racines de 1 dans }k^*.\tag*{(4.2.2)}\]
\scanpage{80}
\mspage{II. 28}
% archive sheet number: 79
% lines 1–4
Alors, si \deleted{\unsure{$\ell\nmid$}} $\ell\nmid N$, \deleted{on aura $\mu_\ell(k^*)$} $\pi_0$ opère non trivialt sur le \deleted{\unsure{groupe}} \inserted{$\mathbb F_\ell$-espace} vectoriel de dim 1 $\mu_\ell(\bar k^*)^{\unsure{\otimes1|\otimes-1}}$, donc sur son dual $V$, donc \deleted{la relation} $V^{\pi_0}=0$, ce qui implique
% line 5
\[H^1(\bar X,\mathbb Z/\ell\mathbb Z)^\pi\subset H^1(\bar X',\mathbb Z/\ell\mathbb Z)^\pi\quad\text{si }\ell\deleted{\unsure{\mid}}\nmid N\tag*{(4.2.3)}\]
% lines 6–8
\deleted{Or on a, si $A$ est la jacobienne de $X'$}
\[H^1(\bar X',\mathbb Z/\ell\mathbb Z)\simeq\deleted{\unsure{\operatorname{Hom}({}_\ell A(\bar k))|T_\ell(A(\bar k))}}{}_{\ell}A(\unsure{\bar k^*|\bar k})^\vee\deleted{\unsure{\otimes\mu_\ell(\bar k)^{\otimes1}}},\]
\deleted{où le $\vee$ désigne le dual à valeurs dans $\mathbb Z/\ell\mathbb Z$.}
% lines 9–13
D'autre part, \deleted{\unsure{pour le} $H^1(\bar X',\mathbb Z/\ell\mathbb Z)^\pi=0$ signifie} comme dans 2.5 on se ramène, pour la relation \deleted{i.e. ${}_\ell A(\bar k)^{\vee\pi}=0$ i.e. $({}_\ell A(\bar k))_\pi=0$,} $H^1(\bar X',\mathbb Z/\ell\mathbb Z)^\pi=0$ pour presque $\unsure{\forall|\text{tout}}\ell$, au cas où $k$ est fini.
% lines 14–18
\deleted{Alors il est bien connu, que} \deleted{\unsure{Soit $f$}} Soit donc \deleted{\unsure{comme}} $f$ le générateur \inserted{de Frobenius} de \deleted{\unsure{car $p($}} $\pi_0\simeq\widehat{\mathbb Z}$, et soit $V={}_\ell A(\bar k)$. Comme $V$ est fini, $\operatorname{Ker}(1-f)_V$ et $\operatorname{Coker}(1-f_V)$ ont même cardinal, donc \deleted{la relation} \deleted{\unsure{est}}
\[({}_\ell A(\bar k))_\pi=0\Longleftrightarrow({}_\ell A(\bar k))^\pi=0\quad\text{i.e. }{}_\ell A(k)=0.\]
% lines 19–23
Or soit $N'$ l'\deleted{nombre d'élé} ordre du groupe \inserted{fini} $A(k)$. On voit que \deleted{\unsure{si $\ell$}}
\[{}_\ell A(k)=0\quad\text{si}\quad\ell\nmid N'.\]
Par suite, on trouve que si $\ell\nmid NN'$, alors
\[H^1(\bar X^{\unsure{\ell|'}},\mathbb Z/\ell\mathbb Z)^\pi=0,\quad\text{qfd.}\]
% lines 24–27; corollary continued in lower margin
\uline{Corollaire} \inserted{4.3} Soit $X$ un schéma \inserted{normal} géométriquement \deleted{intègre} \inserted{conn. de t.f.} sur un corps fini $\mathbb F_{p^{\unsure{\nu|v}}}$. Alors pour tout nb premier $\ell\neq p$, \inserted{l'hom. surj.} $\pi_1(X,\xi)^{\mathrm{ab}}(\ell)\longrightarrow\pi_1(\mathbb F_p)\simeq\widehat{\mathbb Z}$ \deleted{\unsure{$\mathbb Z$ est libre de rang}} a un noyau fini, et pour presque tout $\ell$, c'est un isomorphisme.
\scanpage{81}
\mspage{II. 29}
% archive sheet number: 80
% lines 1–2
5. \uline{Réduction du} \uline{th. de monodromie} \uline{pour} \deleted{\unsure{$H^1$}} \inserted{$H^1(\bar X,\mathbb Z_\ell)$} \uline{au cas} \uline{d'une courbe projective lisse} \uline{géom. connexe}.
% lines 3–5
5.1. La suite spectrale de descente \deleted{\unsure{étale}} \deleted{\unsure{nous}} pour un morphisme fini (SGA 4 \uline{VIII}) nous ramène au cas $X$ \uline{normal}.
% lines 6–9
5.2. \inserted{On peut déjà supposer $X$ affine, et} Une extension \deleted{radicielle} \inserted{finie} du corps de base nous ramène au cas où de plus $X$ est un ouvert d'un schéma projectif \deleted{normal} \inserted{$X'$} géom.$^{\mathrm t}$ connexe et géom.$^{\mathrm t}$ normal.
% lines 10–13
5.3. Dans ce cas, 3.2 (qui implique que
\[H^1(\bar X,\mathbb Z_\ell)-H^1(\bar C,\mathbb Z_\ell)\]
est injectif) nous ramène au cas \deleted{\unsure{de dim}} $X'$ est \inserted{une \uline{courbe}} \deleted{\unsure{géom. normale connexe}} projective et lisse.
% lines 14–17
[\uline{NB} \deleted{Si $K=k(t_1,\ldots,t_r)$, il est} On choisit arbitrairement un anneau de valuation $V'$ dans $K'=K(t)$, \inserted{qui} domine $V$ \unsure{\}}; alors l'image de $I'-I$ est d'indice fini dans $\unsure{V|I}$].
% lines 18–23
Dans ce cas, \deleted{la suite} on peut encore supposer que \inserted{les pts de} $X'-X$ sont rat / $K$. Alors \inserted{dans} la suite exacte (2.5.1), l'opération de $\bar I$ sur les $H^{\unsure{1|2}}_x(\bar X,\mathbb Z_\ell)\simeq\mathbb Z_\ell(-1)$ est triviale, donc il reste à prouver que l'action sur $H^2(\bar X',\mathbb Z_\ell)$ est quasi-unipotente, ce qui est la réduction annoncée.
\scanpage{82}
\mspage{II. 30}
% archive sheet number: 81
% lines 1–6
\uline{Remarque} 5.4. Dans Exp \uline{III}, nous prouverons \unsure{$=$|$\simeq$}, quand $X_K$ est une courbe \deleted{projective} \inserted{séparée} lisse sur $K$, que l'action de $I$ sur $E=H^1(X_{\bar K},\mathbb Z_\ell)$ est \unsure{pot$^{\mathrm t}$|ess$^{\mathrm t}$}-unipotente d'\uline{échelon} 2, i.e. qu'il existe un ss-groupe ouvert \deleted{\unsure{invariant}} \inserted{$U$} de $I$ tel que \deleted{\unsure{$I$}} $U$ opère \deleted{de façon} trivialement sur \deleted{\unsure{$E$}} $E/E^U$.
% lines 7–10
On voit alors facilement, en suivant les réductions 5.1, 5.2, 5.3, \deleted{qu'il en est de même} pour tout \deleted{\unsure{tt}} $X$ \deleted{normal} de type fini sur $K$, l'action de $I$ est \deleted{quasi} ess$^{\mathrm t}$ unipotente d'échelon 3 (et $=$ d'échelon 2 \deleted{\unsure{si}} si $X$ est normal).
% lines 11–18; lower half
\uline{Bibliographie}

[1] A. Grothendieck, Un th. sur les hom. de schémas abéliens, Inv. math 2, 59-78 (1966).

(2) S. Lang,

[3] J.P. Serre--J. Tate, Good reduction of abelian varieties.

[\unsure{4|3}] Shimura--Taniyama, Complex multiplication of abelian varieties, Pub. of Math. Soc. of Japan, n° 6 (1961).
